% GENERATED FILE. Edit canonical lessons, then run npm run generate:books.
% canonical-source-hash: fnv1a64:79afe391ebf30c00
% canonical-lessons: JA-C488-kenchiku, JA-C488-zaimoku, JA-C488-renga, JA-C488-tetsu, JA-C488-kouji

\chapter{Building Work}
\label{ch:ja-building-work}
\hlchaptermodality{488}

\begin{chapteropening}
\textbf{By the end of this chapter:} \emph{I can talk about building work and what buildings are made of.}

Complete the last lesson of chapter 488: I can talk about building work and what buildings are made of.
\end{chapteropening}

\section[\texorpdfstring{\ja{けんちく} (kenchiku)}{けんちく (kenchiku)}]{\ja{けんちく} (kenchiku) — architecture; building}

\label{lesson:JA-C488-kenchiku}



\begin{quote}
Before the new one: say the Japanese for sales, then the Japanese for an employee.
\end{quote}

\subsection*{You'll want to know: \ja{けんちく}}
\textbf{\ja{けんちく}} — \emph{kenchiku} — \textquotedblleft{}architecture; building\textquotedblright{}.

\begin{quote}
\textbf{\ja{あにの} \ja{しごとは} \ja{けんちくです。}} — \emph{ani no shigoto wa kenchiku desu} — \textquotedblleft{}My elder brother works in building.\textquotedblright{}
\end{quote}

\begin{cousinweb}
\textbf{\ja{けんちく}} (\emph{kenchiku}) is Sino-Japanese (kango): two Chinese characters read with their on'yomi, readings borrowed from Chinese. They are \emph{ken}, \textquotedblleft{}build\textquotedblright{} (Middle Chinese \emph{kjonH}) and \emph{chiku}, \textquotedblleft{}fabricate\textquotedblright{} (KANJIDIC2; Middle Chinese in Baxter's transcription).
\end{cousinweb}

\begin{grammarlens}[title={What you've built}]
One more word for describing what you lived and felt.
\end{grammarlens}

\subsection*{Your turn}
\begin{itemize}
\raggedright
  \item \emph{Say it:} \emph{kenchiku}
  \item \emph{Say it:} \emph{kenchiku}, once more
  \item \emph{Say it:} \emph{ani no shigoto wa kenchiku desu} — My elder brother works in building.
  \item \emph{Recall:} say the Japanese for sales, then the Japanese for an employee, then say \emph{kenchiku} again
\end{itemize}

\begin{tcolorbox}[breakable,colback=teal!4,colframe=teal!35!black,title={Before you move on}]
What does \ja{けんちく} mean? (\textquotedblleft{}Architecture; building\textquotedblright{}.) Say it once more.
\end{tcolorbox}
\section[\texorpdfstring{\ja{ざいもく} (zaimoku)}{ざいもく (zaimoku)}]{\ja{ざいもく} (zaimoku) — timber; lumber}

\label{lesson:JA-C488-zaimoku}



\begin{quote}
Before the new one: say the Japanese for an employee, then the Japanese for architecture.
\end{quote}

\subsection*{You'll want to know: \ja{ざいもく}}
\textbf{\ja{ざいもく}} — \emph{zaimoku} — \textquotedblleft{}timber; lumber\textquotedblright{}.

\begin{quote}
\textbf{\ja{くるまで} \ja{ざいもくを} \ja{はこびます。}} — \emph{kuruma de zaimoku o hakobimasu} — \textquotedblleft{}We carry the timber by car.\textquotedblright{}
\end{quote}

\begin{cousinweb}
\textbf{\ja{ざいもく}} (\emph{zaimoku}) is Sino-Japanese (kango): two Chinese characters read with their on'yomi, readings borrowed from Chinese. They are \emph{zai}, \textquotedblleft{}lumber\textquotedblright{} and \emph{moku}, \textquotedblleft{}tree\textquotedblright{} (Middle Chinese \emph{muwk}) (KANJIDIC2; Middle Chinese in Baxter's transcription).
\end{cousinweb}

\begin{grammarlens}[title={What you've built}]
One more word for describing what you lived and felt.
\end{grammarlens}

\subsection*{Your turn}
\begin{itemize}
\raggedright
  \item \emph{Say it:} \emph{zaimoku}
  \item \emph{Say it:} \emph{zaimoku}, once more
  \item \emph{Say it:} \emph{kuruma de zaimoku o hakobimasu} — We carry the timber by car.
  \item \emph{Recall:} say the Japanese for an employee, then the Japanese for architecture, then say \emph{zaimoku} again
\end{itemize}

\begin{tcolorbox}[breakable,colback=teal!4,colframe=teal!35!black,title={Before you move on}]
What does \ja{ざいもく} mean? (\textquotedblleft{}Timber; lumber\textquotedblright{}.) Say it once more.
\end{tcolorbox}
\section[\texorpdfstring{\ja{れんが} (renga)}{れんが (renga)}]{\ja{れんが} (renga) — a brick}

\label{lesson:JA-C488-renga}



\begin{quote}
Before the new one: say the Japanese for architecture, then the Japanese for timber.
\end{quote}

\subsection*{You'll want to know: \ja{れんが}}
\textbf{\ja{れんが}} — \emph{renga} — \textquotedblleft{}a brick\textquotedblright{}.

\begin{quote}
\textbf{\ja{この} \ja{えきは} \ja{れんがの} \ja{たてものです。}} — \emph{kono eki wa renga no tatemono desu} — \textquotedblleft{}This station is a brick building.\textquotedblright{}
\end{quote}

\begin{cousinweb}
\textbf{\ja{れんが}} (\emph{renga}) is Sino-Japanese (kango): two Chinese characters read with their on'yomi, readings borrowed from Chinese. They are \emph{ren}, \textquotedblleft{}refine metals\textquotedblright{} and \emph{ga}, \textquotedblleft{}tile\textquotedblright{} (Middle Chinese \emph{ngwaeX}) (KANJIDIC2; Middle Chinese in Baxter's transcription).
\end{cousinweb}

\begin{grammarlens}[title={What you've built}]
One more word for describing what you lived and felt.
\end{grammarlens}

\subsection*{Your turn}
\begin{itemize}
\raggedright
  \item \emph{Say it:} \emph{renga}
  \item \emph{Say it:} \emph{renga}, once more
  \item \emph{Say it:} \emph{kono eki wa renga no tatemono desu} — This station is a brick building.
  \item \emph{Recall:} say the Japanese for architecture, then the Japanese for timber, then say \emph{renga} again
\end{itemize}

\begin{tcolorbox}[breakable,colback=teal!4,colframe=teal!35!black,title={Before you move on}]
What does \ja{れんが} mean? (\textquotedblleft{}A brick\textquotedblright{}.) Say it once more.
\end{tcolorbox}
\section[\texorpdfstring{\ja{てつ} (tetsu)}{てつ (tetsu)}]{\ja{てつ} (tetsu) — iron}

\label{lesson:JA-C488-tetsu}



\begin{quote}
Before the new one: say the Japanese for timber, then the Japanese for a brick.
\end{quote}

\subsection*{You'll want to know: \ja{てつ}}
\textbf{\ja{てつ}} — \emph{tetsu} — \textquotedblleft{}iron\textquotedblright{}.

\begin{quote}
\textbf{\ja{てつの} \ja{なべは} \ja{おもいです。}} — \emph{tetsu no nabe wa omoi desu} — \textquotedblleft{}An iron pot is heavy.\textquotedblright{}
\end{quote}

\begin{cousinweb}
\textbf{\ja{てつ}} (\emph{tetsu}) is Sino-Japanese (kango): one Chinese character read with its on'yomi, reading borrowed from Chinese. It is \emph{tetsu}, \textquotedblleft{}iron\textquotedblright{} (Middle Chinese \emph{thet}) (KANJIDIC2; Middle Chinese in Baxter's transcription).
\end{cousinweb}

\begin{grammarlens}[title={What you've built}]
One more word for describing what you lived and felt.
\end{grammarlens}

\subsection*{Your turn}
\begin{itemize}
\raggedright
  \item \emph{Say it:} \emph{tetsu}
  \item \emph{Say it:} \emph{tetsu}, once more
  \item \emph{Say it:} \emph{tetsu no nabe wa omoi desu} — An iron pot is heavy.
  \item \emph{Recall:} say the Japanese for timber, then the Japanese for a brick, then say \emph{tetsu} again
\end{itemize}

\begin{tcolorbox}[breakable,colback=teal!4,colframe=teal!35!black,title={Before you move on}]
What does \ja{てつ} mean? (\textquotedblleft{}Iron\textquotedblright{}.) Say it once more.
\end{tcolorbox}
\section[\texorpdfstring{\ja{こうじ} (kōji)}{こうじ (kōji)}]{\ja{こうじ} (kōji) — construction work}

\label{lesson:JA-C488-kouji}



\begin{quote}
Before the new one: say the Japanese for a brick, then the Japanese for iron.
\end{quote}

\subsection*{You'll want to know: \ja{こうじ}}
\textbf{\ja{こうじ}} — \emph{kōji} — \textquotedblleft{}construction work\textquotedblright{}.

\begin{quote}
\textbf{\ja{みちで} \ja{こうじを} \ja{しています。}} — \emph{michi de kōji o shiteimasu} — \textquotedblleft{}There is construction work on the road.\textquotedblright{}
\end{quote}

\begin{cousinweb}
\textbf{\ja{こうじ}} (\emph{kōji}) is Sino-Japanese (kango): two Chinese characters read with their on'yomi, readings borrowed from Chinese. They are \emph{kō}, \textquotedblleft{}construction\textquotedblright{} (Middle Chinese \emph{kuwng}) and \emph{ji}, \textquotedblleft{}matter\textquotedblright{} (\emph{dzriH}) (KANJIDIC2; Middle Chinese in Baxter's transcription).
\end{cousinweb}

\begin{grammarlens}[title={What you've built}]
One more word for describing what you lived and felt.
\end{grammarlens}

\subsection*{Your turn}
\begin{itemize}
\raggedright
  \item \emph{Say it:} \emph{kōji}
  \item \emph{Say it:} \emph{kōji}, once more
  \item \emph{Say it:} \emph{michi de kōji o shiteimasu} — There is construction work on the road.
  \item \emph{Recall:} say the Japanese for a brick, then the Japanese for iron, then say \emph{kōji} again
\end{itemize}

\begin{tcolorbox}[breakable,colback=teal!4,colframe=teal!35!black,title={Before you move on}]
What does \ja{こうじ} mean? (\textquotedblleft{}Construction work\textquotedblright{}.) Say it once more.
\end{tcolorbox}
